\documentclass[11pt,UTF8,fontset=fandol]{ctexart}
\usepackage{amsmath,amssymb,amsfonts}
\usepackage{graphicx,float,caption,subcaption}
\usepackage{booktabs,array,multirow,makecell}
\usepackage{geometry,enumitem,xcolor}
\usepackage[hidelinks]{hyperref}
\usepackage{url}
\geometry{margin=2.3cm}
\graphicspath{{/Users/yye/project/论文/assets/eto/}}
\captionsetup{font=small,labelfont=bf}
\setlength{\parskip}{0.4em}

% 小白知识卡 / 重点框 的简单样式
\newcommand{\zhcard}[2]{%
  \par\noindent\textbf{#1}\quad #2\par}

\title{\textbf{试错学习：面向 LLM 智能体的\\基于探索的轨迹优化（ETO）}}
\author{Yifan Song, Da Yin, Xiang Yue, Jie Huang, Sujian Li, Bill Yuchen Lin\\
\small Allen Institute for AI \quad 北京大学（Peking University）\quad UCLA \quad Ohio State University \quad UIUC}
\date{}

\begin{document}
\maketitle

\begin{center}
\fbox{\parbox{0.95\linewidth}{\small
\textbf{英文原标题：} Trial and Error: Exploration-Based Trajectory Optimization for LLM Agents \\
\textbf{发表：} ACL 2024 主会（Main Conference）\quad\textbf{引用数：} 约 181（Semantic Scholar，2026-06-07）\\
\textbf{代码与数据：} \url{https://github.com/Yifan-Song793/ETO} \\
\textbf{一句话定位：} 让开源大模型智能体\textbf{从自己的失败里学习}——先模仿专家学会基本操作，再到环境里真刀真枪地试、收集失败案例，把“失败轨迹”和“成功轨迹”配成对，用 DPO 这种对比学习的方式“奖善惩失”，并可反复迭代越练越强。
}}
\end{center}

\medskip

% =====================================================================
\section*{小白知识卡：读这篇前先搞懂这几个词}
\addcontentsline{toc}{section}{小白知识卡}

\zhcard{智能体（agent）}{这里指一个能“看环境—想一步—做动作—再看反馈”循环操作的大模型程序。比如让它去网购、做虚拟科学实验、在虚拟房间里做家务，它要一步步执行直到完成任务。}

\zhcard{轨迹（trajectory）}{智能体完成一个任务从头到尾的完整记录：任务指令 + 第 1 个动作 + 环境反馈 + 第 2 个动作 + ……一直到结束。一条轨迹就是“这一局是怎么打的”的全过程。}

\zhcard{奖励（reward）}{环境给这条轨迹打的分，一般是 0 到 1。1 表示任务完美完成，0 表示失败。它是判断“成功还是失败”的唯一依据。}

\zhcard{模仿学习 / 行为克隆（imitation learning / behavioral cloning，BC）}{最朴素的训练法：给模型看一堆“专家是怎么做的”示范，让它照着学、把专家的动作背下来。本文里就是对专家轨迹做\textbf{监督微调（supervised fine-tuning，SFT）}。缺点：只会模仿，没真正去环境里探索过，遇到没见过的情况容易翻车。}

\zhcard{ReAct 与思维链（Chain-of-Thought，CoT）}{ReAct 是一种让智能体“边想边做”的格式：每做一个动作前，先用自然语言写一段推理（也就是 CoT），再给出动作。本文所有智能体都用这种“先想后做”的格式。}

\zhcard{强化学习（reinforcement learning，RL）与 PPO}{让智能体自己在环境里试、根据拿到的奖励来调整策略的训练范式。PPO（Proximal Policy Optimization，近端策略优化）是其中最常用的在线 RL 算法。它能真正“探索”，但在大模型上\textbf{训练不稳定、效率低}，很难调。}

\zhcard{DPO（Direct Preference Optimization，直接偏好优化）}{一种“绕开 RL 复杂流程”的训练法。你只要给它一对数据“A 比 B 好”，它就能直接调模型：\textbf{抬高好样本的概率、压低坏样本的概率}。本文的核心武器：把“成功轨迹 vs 失败轨迹”当成这样一对偏好数据喂给 DPO。}

\zhcard{信用分配（credit assignment）}{一条轨迹好几十步，最后只给一个总分——到底是哪一步做对/做错导致了这个结果？把功劳或锅分到具体某一步，就叫信用分配。这是序列决策里出了名的难题，本文末尾也坦承自己在这点上做了简化。}

\zhcard{分布外泛化（out-of-distribution，OOD / unseen）}{测试时遇到训练时\textbf{没见过的新任务变体}。比如训练时学的是“烧开水”，测试时却要“熔化铅”。能不能在这种新场景上也做好，衡量的就是泛化能力。}

% =====================================================================
\section{摘要全译}

大语言模型（Large Language Models，LLMs）已经成为各类自主智能体（autonomous agent）系统不可或缺的组成部分。在本研究中，我们提出一种基于探索的轨迹优化方法，称为 ETO（Exploration-based Trajectory Optimization）。这一学习方法旨在提升开源大模型智能体的性能。与以往\textbf{只在成功的专家轨迹上训练}的研究不同，我们的方法让智能体能够从\textbf{自身探索的失败}中学习。借助一个迭代式的优化框架，这带来了性能上的提升。在探索阶段，智能体一边完成给定任务、一边与环境交互，收集失败轨迹，从而构造出成对的“对比轨迹”。在随后的训练阶段，智能体利用这些轨迹偏好对，采用 DPO（Rafailov et al., 2023）等对比学习方法来更新自己的策略。这种“探索—训练”的迭代循环促使智能体持续进步。我们在三项复杂任务上的实验表明，ETO 始终大幅超越基线方法的表现。此外，对“任务求解效率”以及“缺少专家轨迹场景下的潜力”的考察，进一步印证了本方法的有效性。

% =====================================================================
\section{引言：为什么要让智能体“吃一堑长一智”}

\textbf{背景痛点。} 大模型在网页浏览、具身任务（embodied tasks，指在模拟的物理/家居环境里操作）、多模态推理、复杂问答等交互任务上展现了惊人能力。但用 ChatGPT、GPT-4 当“大脑”毕竟成本高、不可控；而\textbf{开源大模型当智能体时，效果比 GPT-4 差很多}（Liu et al., 2023 等）。所以业界希望把开源模型训练成好用的智能体。

\textbf{现有做法的局限。} 主流做法是\textbf{模仿学习}：拿专家轨迹对开源模型做 SFT（如 FireAct、AgentTuning、Lumos 等）。这相当于让模型“背标准答案”。问题在于：它\textbf{完全依赖专家示范，从没真正去目标环境里探索过}，于是学到的往往是次优策略，泛化能力也受限——遇到稍有变化的新任务就容易出错。

\textbf{本文的灵感与思路。} 人类学习不仅靠看成功示范，更靠\textbf{在试错（trial and error）中亲历失败}。受此启发，作者提出 ETO：不再只用成功轨迹，而是\textbf{利用当前策略探索时产生的失败}来强化学习过程。整套流程是（见图~\ref{fig:overview}）：

\begin{enumerate}[leftmargin=2em]
\item 先用 SFT 行为克隆训练出一个\textbf{基础智能体}（base agent）；
\item \textbf{探索阶段：} 让这个基础智能体去目标环境里实际跑一遍训练任务，收集它\textbf{失败的轨迹}，并和之前已有的\textbf{专家成功轨迹}配成对；
\item \textbf{训练阶段：} 用 DPO 损失（Rafailov et al., 2023），拿这些“失败 vs 成功”的对比轨迹对去微调模型，让它\textbf{偏好成功、厌恶失败}；
\item 上述过程可\textbf{多轮迭代}：用上一轮练出来的智能体再去探索、再收集新失败、再训练，循环往复。
\end{enumerate}

\textbf{实验与结论速览。} 作者在三个代表性数据集上验证：WebShop（网页购物导航）、ScienceWorld（模拟科学实验）、ALFWorld（具身家务）。结果是：ETO 始终大幅超过 SFT 行为克隆和其他强基线；在 ScienceWorld 的分布外（OOD）测试集上比 SFT 提升达 22\%，泛化能力突出；它还能用更少的动作步数拿到更高奖励（效率更高）；甚至在\textbf{完全没有专家轨迹}的极端情形下，靠“自我对弈（self-play）”也能有不错表现。

\textbf{三点贡献。} （1）\textbf{方法}：提出 ETO，一种迭代收集失败轨迹、用对比学习精炼策略的算法；（2）\textbf{评测}：在三项复杂交互任务上全面实验，超越 SFT 与其他强基线；（3）\textbf{分析}：从分布外泛化、动作效率、以及“无需专家轨迹的可行性”等多个角度深入验证有效性。

% =====================================================================
\section{任务形式化：把“做任务”写成数学语言}

作者把“带环境反馈的智能体任务”形式化为一个\textbf{部分可观测马尔可夫决策过程}（partially observable Markov decision process，POMDP），记作六元组 $(\mathcal{U}, \mathcal{S}, \mathcal{A}, \mathcal{O}, \mathcal{T}, \mathcal{R})$：

\begin{itemize}[leftmargin=2em,itemsep=2pt]
\item $\mathcal{U}$：指令空间（任务怎么描述）；$\mathcal{S}$：状态空间（环境真实状态）；$\mathcal{A}$：动作空间；$\mathcal{O}$：观测空间（智能体能看到的反馈）；
\item $\mathcal{T}: \mathcal{S}\times\mathcal{A}\to\mathcal{S}$：转移函数（做了动作后状态怎么变）；$\mathcal{R}: \mathcal{S}\times\mathcal{A}\to[0,1]$：奖励函数。
\item \textbf{关键点}：在大模型智能体场景里，$\mathcal{U}, \mathcal{A}, \mathcal{O}$ 都是\textbf{自然语言}的子集——也就是说，指令、动作、观测都是文字。
\end{itemize}

给定一条任务指令 $u\in\mathcal{U}$，参数为 $\theta$ 的大模型智能体按其策略 $\pi_\theta$ 生成第一个动作 $a_1\sim\pi_\theta(\cdot|u)$。这个动作会改变隐状态 $s_t$、并返回一个观测 $o_t$；接着在第 $t{+}1$ 步生成 $a_{t+1}\sim\pi_\theta(\cdot|u, a_1, o_1, \dots, o_{t-1}, a_t)$。如此循环，直到任务完成或超过最大步数。整条轨迹记为：

\begin{equation*}
e = (u, a_1, o_1, \dots, o_{n-1}, a_n) \sim \pi_\theta(e|u), \tag{1}
\end{equation*}
\begin{equation*}
\pi_\theta(e|u) = \prod_{j=1}^{n} \pi_\theta(a_j | u, a_1, o_1, \dots, o_{j-1}), \tag{2}
\end{equation*}

其中 $n$ 是轨迹长度。\textbf{公式 (2) 拆解（很重要）}：一整条轨迹的概率，就等于\textbf{逐步动作概率连乘}——第 $j$ 步动作 $a_j$ 是在“看到此前所有动作和观测”的条件下生成的。注意连乘里只算\textbf{动作} $a_j$ 的概率，不算环境给的观测 $o$（因为观测不是模型生成的）。最后环境算出总奖励 $r(u,e)\in[0,1]$，等于 1 即任务成功。

% =====================================================================
\section{方法：ETO 怎么从失败里学习}

ETO 先用行为克隆造一个基础智能体，然后在此基础上\textbf{迭代式地“从试错中精炼策略”}。

\subsection{第一步：行为克隆（打底）}

行为克隆（BC）就是在专家交互轨迹上做 SFT，它是搭建强智能体的\textbf{可靠起点}。本文用 ReAct 风格的轨迹来做 BC——即每个动作前都额外生成一段 CoT 推理。为简化记号，下文用 $a$ 同时代表“CoT + 动作”这一整块。

给定专家轨迹数据集 $D = \{(u, e)^{(i)}\}_{i=1}^{|D|}$，在自回归损失上微调出基础智能体 $\pi_{\text{base}}$：

\begin{equation*}
\mathcal{L}_{\text{SFT}}(\pi_\theta) = -\mathbb{E}_{e\sim D}\left[\pi_\theta(e|u)\right], \tag{3}
\end{equation*}

其中 $e=(u, a_1, o_1, \dots, o_{n-1}, a_n)\sim D$ 是一条专家轨迹。具体实现时，把指令、动作、观测拼成一段文本序列 $t=\text{concat}(u, a_1, o_1, \dots, o_{n-1}, a_n) = (t_1, t_2, \dots, t_l)$，$t_k$ 是第 $k$ 个 token。于是轨迹概率可由下式得到（\textbf{把任务描述和观测部分的 token 都“遮住”不算损失}）：

\begin{equation*}
\pi_\theta(e|u) = -\sum_{k} \log \pi_\theta(t_k | t_{<k}) \times \mathbb{1}(t_k \in \mathcal{A}), \tag{5}
\end{equation*}

\textbf{公式 (5) 拆解}：$\mathbb{1}(t_k\in\mathcal{A})$ 是一个指示函数——只有当这个 token 属于\textbf{智能体自己产出的动作}时才取 1，否则取 0。换句话说，训练时\textbf{只对“模型该说的话（动作/CoT）”计算损失，对指令和环境观测不计损失}。这很合理：环境反馈不是模型该背的，强行学反而有害。

\subsection{第二步：从探索失败中学习（核心）}

\textbf{为什么不直接上 RL？} 行为克隆只会模仿专家、不会探索，策略往往次优。要练得更强，理应让模型也去探索失败轨迹——最直接的办法就是强化学习，其目标是：

\begin{equation*}
\max_{\pi_\theta}\; \mathbb{E}_{u\sim D,\, e\sim\pi_\theta(e|u)}\left[r(u,e)\right] - \beta\, D_{\text{KL}}\left[\pi_\theta(e|u)\,\|\,\pi_{\text{ref}}(e|u)\right], \tag{6}
\end{equation*}

\textbf{公式 (6) 拆解}：前一项是\textbf{最大化期望奖励}（多拿高分），后一项是一个带权重 $\beta$ 的 KL 散度约束，作用是\textbf{别让新策略 $\pi_\theta$ 偏离参考策略 $\pi_{\text{ref}}$（也就是基础智能体 $\pi_{\text{base}}$）太远}，防止训崩。实践中待训练的 $\pi_\theta$ 也用 $\pi_{\text{base}}$ 初始化。这个目标可以用 PPO 等 RL 方法来解。

\textbf{但是——} 在大模型智能体上直接做\textbf{在线 RL 有现实困难：不稳定、效率低}。因此作者改走\textbf{迭代式离线学习}路线：用\textbf{对比轨迹数据}来训练，把整个过程组织成“探索—训练”循环（见图~\ref{fig:overview}）。

\paragraph{探索阶段（Exploration Phase）。} 基础智能体 $\pi_{\text{base}}$ 在 BC 用过的那批训练指令上去环境里跑一遍，得到自己的轨迹：

\begin{equation*}
\hat{e} = (u, \hat{a}_1, \hat{o}_1, \dots, \hat{o}_{m-1}, \hat{a}_m) \sim \pi_{\text{base}}(e|u). \tag{7}
\end{equation*}

环境为这条轨迹 $\hat{e}$ 返回奖励 $\hat{r}$。然后构造\textbf{“失败—成功”轨迹对}，记作 $e_w \succ e_l \mid u$：$e_w$ 是高奖励（胜出）轨迹、$e_l$ 是低奖励（落败）轨迹，分别从\textbf{专家轨迹 $e$} 和\textbf{智能体自己跑出的轨迹 $\hat{e}$} 中按奖励高低选出。\textbf{注意两条规则}：（1）只收集\textbf{两条轨迹奖励不同}的对；（2）如果智能体和专家\textbf{都成功了}，这一对就直接\textbf{丢弃}（没有对比价值）。最终得到对比轨迹数据集 $D_p = \{(u, e_w, e_l)^{(i)}\}_{i=1}^{|D_p|}$。

\paragraph{训练阶段（Training Phase）。} 用轨迹对里的“成败对比信息”来更新策略。给定 $e_w \succ e_l \mid u$，这种成败关系用 \textbf{Bradley-Terry（BT）模型}建模：

\begin{equation*}
p(e_w \succ e_l \mid u) = \frac{\exp(r(u, e_w))}{\exp(r(u, e_w)) + \exp(r(u, e_l))}. \tag{8}
\end{equation*}

\textbf{公式 (8) 直觉}：奖励越高的轨迹，“它更优”的概率就越大——这就是个用奖励差异决定胜负概率的 softmax。再借助 DPO 的关键推导：在目标 (6) 的最优策略 $\pi_r$ 下，奖励可写成 $r(u,e)=\beta\log\frac{\pi_r(e|u)}{\pi_{\text{ref}}(e|u)} + \beta\log Z(x)$（其中 $Z(u)$ 是配分函数）。把它代回 (8)，配分函数恰好抵消，得到\textbf{直接用策略表示的 BT 模型}，从而推出 DPO 损失：

\begin{equation*}
\mathcal{L}_{\text{DPO}}(\pi_\theta; \pi_{\text{ref}}) = -\,\mathbb{E}_{(u,e_w,e_l)\sim D_p}\left[\log\sigma\!\left(\beta\log\frac{\pi_\theta(e_w|u)}{\pi_\theta(e_l|u)} - \beta\log\frac{\pi_{\text{ref}}(e_w|u)}{\pi_{\text{ref}}(e_l|u)}\right)\right]. \tag{11}
\end{equation*}

\textbf{公式 (11) 拆解（全文最关键的一条）}：$\sigma$ 是 sigmoid 函数。括号里在比较两件事——“当前策略 $\pi_\theta$ 对成功轨迹 vs 失败轨迹的相对偏好”，减去“参考策略 $\pi_{\text{ref}}$ 的同样偏好”。优化它的效果就是：\textbf{抬高成功轨迹 $e_w$ 的概率、压低失败轨迹 $e_l$ 的概率}，同时用参考策略那一项当“锚”，约束模型别跑偏。妙处在于：\textbf{这其实是 RL 目标 (6) 的等价改写}，它在直接最大化最终奖励，却\textbf{完全不用做 RL 那套不稳定的优化}。

\subsection{第三步：迭代（越练越强）}

训练完一轮后，更新过的策略可以\textbf{再去收集新的失败案例}、构造新的对比对，再训练一轮。完整流程见算法 1（Algorithm 1）：先做 $T_1$ 步行为克隆得到基础智能体；然后做 $I$ 轮迭代，每轮里令 $\pi_{\text{base}}=\pi_{\text{ref}}=\pi_\theta$，跑出智能体轨迹、与专家轨迹比奖励配成对、构造 $D_p$，再做 $T_2$ 步 DPO 训练。如此循环输出最终策略 $\pi_\theta$。

% =====================================================================
\section{方法图解：一张图看懂数据怎么流动}

\begin{figure}[H]\centering
  \includegraphics[width=0.98\linewidth]{cand_p03_k00.pdf}
  \caption{ETO 总览（原文图 2）。从一个经行为克隆训练得到的基础大模型智能体出发，本方法让智能体迭代地收集失败轨迹，并通过不断从“失败—成功”对比轨迹对中学习来更新策略。}
  \label{fig:overview}
\end{figure}

\textbf{图里在说什么、数据怎么流动？} 整张图从左到右是三个阶段，用一个具体任务“把两部手机放到桌上（put two cellphones in desk）”串起来：

\begin{itemize}[leftmargin=2em,itemsep=3pt]
\item \textbf{左：行为克隆（Behavioral Clone）。} 拿一条专家成功轨迹（Expert Trajectory，最后 Reward=1.0，带金牌图标）对模型做 SFT，得到一个\textbf{会基本操作}的基础智能体（图中是只羊驼/Llama 的形象）。
\item \textbf{中：探索阶段（Exploration Phase）。} 让这个基础智能体（紫色机器人）\textbf{真的去环境里跑这些任务}，把它\textbf{失败的轨迹收集起来（Collect Failures）}。图中下方一摞屏幕就代表它跑出的各种（多半失败的）尝试。
\item \textbf{右：训练阶段（Training Phase）。} 把\textbf{失败轨迹（左列，Reward=0.5，红叉）}和\textbf{成功轨迹（右列，Reward=1.0，绿勾）}\textbf{并排配成一对}，喂给 DPO 做对比学习（“Learning from Contrastive Failure-Success Trajectories”）。对照两列能直观看出差异：失败的那次先去了 table 1/table 2 都没找到第二部手机，最后只拿 0.5 分；成功的那次去了 shelf 1 拿到手机、顺利放进桌子，拿满 1.0 分。模型就是从\textbf{这种“同一个任务、一成一败”的对照}里学到该怎么做、不该怎么做。
\item \textbf{回路箭头（中间的循环图标）}：训练完的新策略可以\textbf{再回到探索阶段}收集新失败，形成迭代闭环。
\end{itemize}

一句话总结这张图：\textbf{先模仿打底 → 放出去试错收集失败 → 把成败配对用 DPO 对比学习 → 循环加强}。

% =====================================================================
\section{实验解读：数字说明了什么}

\subsection{用了什么数据、比了哪些基线}

\textbf{三个交互环境（数据集见表~\ref{tab:data}）：}
\begin{itemize}[leftmargin=2em,itemsep=2pt]
\item \textbf{WebShop}：网购环境，按指令在购物网站上搜索、点选、下单，奖励按商品属性和价格的匹配度给（0$\sim$1 连续值）。
\item \textbf{ScienceWorld}：文字版虚拟环境，做小学级科学实验（热力学、电路等 10 类任务），每个任务有若干子目标，奖励按完成的子目标算。\textbf{测试集是“关键未见变体”}（如训练学“烧开水”、测试要“熔化铅”），专门考泛化。
\item \textbf{ALFWorld}：具身家务环境，完成“把某物放到某处”这类指令；\textbf{只给二元奖励}（成功=1/失败=0），没有中间分。它分 seen（同分布）和 unseen（分布外）两个测试集。
\end{itemize}

\begin{table}[H]\centering
\caption{数据集统计（原文表 1）。“Test-Seen / Test-Unseen”分别是“见过场景 / 未见场景”的测试集；“\#Turns”是专家轨迹的平均交互轮数。}
\label{tab:data}
\begin{tabular}{lcccc}
\toprule
数据集 & \#Train & \#Test-Seen & \#Test-Unseen & \#Turns \\
\midrule
WebShop      & 1938 & 200 & --  & 4.9 \\
ScienceWorld & 1483 & 194 & 241 & 14.4 \\
ALFWorld     & 3321 & 140 & 134 & 10.1 \\
\bottomrule
\end{tabular}
\end{table}

\textbf{对比的基线方法：}（基座主要用 Llama-2-7B-Chat）
\begin{itemize}[leftmargin=2em,itemsep=2pt]
\item \textbf{SFT}：纯行为克隆，也是 ETO 和其他方法的“基础智能体”。
\item \textbf{Best-of-N}：用 SFT 智能体采样 $N=10$ 条，挑奖励最高那条。
\item \textbf{RFT（Rejection sampling Fine-Tuning，拒绝采样微调）}：一个强基线——把\textbf{成功的}采样轨迹加进专家数据集再训练。注意它\textbf{只学成功轨迹}。
\item \textbf{PPO}：直接用 RL 最大化最终奖励的在线方法。
\item 还对比了 GPT-3.5-Turbo、GPT-4、以及没微调过的 Llama-2-7B-Chat。
\end{itemize}

\textbf{评测设置}：统一用 ReAct 格式（先 CoT 再动作），加 1-shot 示例，解码温度设 0.0（确定性输出，Best-of-N 除外）。主指标是\textbf{平均奖励（Average Reward）}。训练用 8 张 A100-80G；ETO 迭代轮数 WebShop/ScienceWorld 为 2、ALFWorld 为 1；DPO 的 $\beta$ 在 WebShop/ScienceWorld 取 0.1、ALFWorld 取 0.5。

\subsection{主结果：ETO 全面领先}

\begin{table}[H]\centering
\caption{各方法在三个数据集上的\textbf{平均奖励}（原文表 2）。“Seen”是任务类型在训练中见过的留出测试集；“Unseen”是含关键未见任务变体的测试集。最优结果加粗。}
\label{tab:main}
\begin{tabular}{lccccc}
\toprule
\multirow{2}{*}{方法} & \multirow{2}{*}{WebShop} & \multicolumn{2}{c}{ScienceWorld} & \multicolumn{2}{c}{ALFWorld} \\
\cmidrule(lr){3-4}\cmidrule(lr){5-6}
 & & Seen & Unseen & Seen & Unseen \\
\midrule
GPT-4                       & 63.2 & 64.8 & 64.4 & 42.9 & 38.1 \\
GPT-3.5-Turbo               & 62.4 & 16.5 & 13.0 & 7.9  & 10.5 \\
Llama-2-7B-Chat             & 17.9 & 3.8  & 3.1  & 0.0  & 0.0  \\
Llama-2-7B-Chat + SFT       & 63.1 & 67.4 & 53.0 & 60.0 & 67.2 \\
Llama-2-7B-Chat + Best-of-N & 63.8 & 70.2 & 57.6 & 62.1 & 69.4 \\
Llama-2-7B-Chat + RFT       & 63.6 & 71.6 & 54.3 & 62.9 & 66.4 \\
Llama-2-7B-Chat + PPO       & 64.2 & 59.4 & 51.7 & 22.1 & 29.1 \\
\textbf{Llama-2-7B-Chat + ETO (ours)} & \textbf{67.4} & \textbf{73.8} & \textbf{65.0} & \textbf{68.6} & \textbf{72.4} \\
\bottomrule
\end{tabular}
\end{table}

\textbf{这些数字说明了什么？}
\begin{itemize}[leftmargin=2em,itemsep=3pt]
\item \textbf{相比 SFT 大幅提升}：WebShop 平均奖励涨约 8\%、ScienceWorld 涨约 9.5\%；在所有数据集上\textbf{超过全部基线}。在 WebShop 上 ETO（67.4）\textbf{甚至超过 GPT-4}（63.2）——一个 7B 开源模型在该任务上压过 GPT-4。
\item \textbf{泛化能力最突出}：在 ScienceWorld-Unseen 上比 SF（53.0$\to$65.0）相对提升约 20\%。说明“从失败中学”确实能提升对未见场景的适应力。
\item \textbf{“只学成功”不够}：RFT 虽比 SFT 好，但因为它只是“成功轨迹的增强版”、\textbf{从不学失败}，提升有限。这反证了\textbf{“失败 vs 成功”的对比信息才是关键}。
\item \textbf{在线 RL（PPO）不靠谱}：PPO 在 WebShop 略有提升，但在另两个环境上很差，尤其 ALFWorld（只有 22.1/29.1）——因为 ALFWorld 只有\textbf{二元奖励}，加上 RL 本身不稳定，PPO 很难学好。这正是 ETO 用离线 DPO 取代在线 RL 的理由。
\end{itemize}

\subsection{换个基座模型也成立}

\begin{table}[H]\centering
\caption{不同基座大模型上的平均奖励（原文表 3）。}
\label{tab:llms}
\begin{tabular}{llccc}
\toprule
基座模型 & 方法 & WebShop & \makecell{SciWorld\\Seen} & \makecell{SciWorld\\Unseen} \\
\midrule
\multirow{2}{*}{Llama-2-13B} & SFT & 66.3 & 68.1 & 57.6 \\
                             & \textbf{ETO} & \textbf{70.7} & \textbf{71.4} & \textbf{68.6} \\
\midrule
\multirow{2}{*}{Mistral-7B}  & SFT & 60.1 & 63.8 & 52.2 \\
                             & \textbf{ETO} & \textbf{66.2} & \textbf{68.5} & \textbf{62.5} \\
\bottomrule
\end{tabular}
\end{table}

换成 Llama-2-13B 和 Mistral-7B，ETO 依然稳定优于 SFT。两个有趣发现：（1）13B 的相对增益比 7B \textbf{小}，暗示 ETO \textbf{对越弱的智能体帮助越大}；（2）Mistral-7B 虽是比 Llama-2-13B 更强的基础语言模型，但当智能体后却\textbf{不如 Llama-2-7B}——说明“基础语言能力强”和“当智能体强”之间\textbf{没有强相关}。

\subsection{效率、迭代次数、对比数据粒度}

\textbf{效率（图 3）}：在 ScienceWorld 里看“多少步拿到多少分”，ETO 能比 SFT \textbf{用更少动作步达到更高奖励}；个别任务（如 15-90、19-23）甚至比 oracle（标准答案智能体）更早拿到满分 100。结论：从失败中学不仅准，还\textbf{更高效}。

\textbf{迭代次数（图 4）}：WebShop 和 ScienceWorld 上，\textbf{前两轮迭代}带来提升，但\textbf{第三轮起开始下降}；ALFWorld 只有第一轮有用，二三轮甚至跌破 SFT。原因：专家轨迹集是\textbf{固定的}、探索又在\textbf{同一批训练集}上做，导致对比数据的\textbf{多样性和数量受限}，多轮后会\textbf{过拟合}；ALFWorld 的\textbf{粗粒度二元奖励}更是雪上加霜。作者建议：未来可用 GPT-4 动态构造更多样的对比数据。

\textbf{对比数据的粒度（表~\ref{tab:gran}）}：作者比较了三种构造对比数据的方式。

\begin{table}[H]\centering
\caption{在 WebShop 上、用不同粒度对比数据训练的智能体平均奖励（原文表 4）。}
\label{tab:gran}
\begin{tabular}{llccc}
\toprule
方法 & 粒度 & lr & $\beta$ & 平均奖励 \\
\midrule
SFT & -- & -- & -- & 63.1 \\
\midrule
\multirow{4}{*}{ETO} & 轨迹级（Trajectory） & 1e-6 & 0.1 & \textbf{67.4} \\
                     & 步级（Step）        & 1e-6 & 0.1 & 8.3 \\
                     & 步级（Step）        & 1e-7 & 0.5 & 62.8 \\
                     & 混合（Mixture）     & 1e-6 & 0.1 & 64.3 \\
\bottomrule
\end{tabular}
\end{table}

\textbf{解读}：\textbf{轨迹级对比}（整条对整条）效果最好（67.4）。\textbf{步级对比}（比较单步“好动作 vs 坏动作”，见附录 C 公式 13）很\textbf{不稳定}——需要更小学习率、更大 $\beta$ 才不崩，且仍只有 62.8。原因是作者只用\textbf{最终奖励}去判断某一步动作的好坏，这种对“单步质量”的估计\textbf{不准确}（这正是前面知识卡里说的“信用分配”难题）。混合策略也不如纯轨迹级。

\subsection{没有专家轨迹怎么办：自我对弈}

\begin{table}[H]\centering
\caption{无行为克隆（不依赖专家轨迹）的自我对弈性能（原文表 5）。上半部分基于 BC 基础智能体；下半部分\textbf{直接用未微调的 Llama} 当起点。$\dagger$ 表示直接提示未微调的 Llama-2-7B-Chat。}
\label{tab:selfplay}
\begin{tabular}{lcc}
\toprule
方法 & 用了 BC？ & 平均奖励 \\
\midrule
SFT & 是 & 63.1 \\
RFT & 是 & 63.6 \\
ETO & 是 & 67.4 \\
\midrule
Llama-2-7B-Chat$^\dagger$ & 否 & 17.9 \\
RFT     & 否 & 48.4 \\
ETO     & 否 & 12.5 \\
RFT+ETO & 否 & 51.2 \\
\bottomrule
\end{tabular}
\end{table}

在\textbf{完全没有专家轨迹}的极端设定下（去掉 BC 阶段、让模型用温度 1.0 自己探索、按奖励把自己的不同轨迹配成偏好对）：单用 ETO \textbf{反而不行}（12.5）——因为没有 BC 打底，模型连基本操作都不会，自己跑出来的轨迹质量太差。而 RFT 能把 17.9 拉到 48.4。\textbf{关键发现}：\textbf{先用 RFT 打底、再上 ETO（RFT+ETO=51.2）}效果最好。启示：没有专家数据时，应\textbf{先想办法让智能体“能跑出一些成功案例”（如 RFT），再让它从失败中对比学习}。

\subsection{ETO 与各类方法的定位（表 6）}

作者用一张表总结 ETO 相对于其他范式的优势：

\begin{table}[H]\centering
\caption{ETO 与各类替代方案的对比（原文表 6）。$\checkmark$ 表示具备该特性、$\times$ 表示不具备。}
\label{tab:vs}
\begin{tabular}{lcccc}
\toprule
方法 & 探索失败 & 奖励驱动 & 效率 & 鲁棒性 \\
\midrule
SFT          & $\times$    & $\times$    & $\checkmark$ & $\checkmark$ \\
RFT          & $\checkmark$ & $\times$    & $\checkmark$ & $\checkmark$ \\
在线 RL（Online RL） & $\checkmark$ & $\checkmark$ & $\times$    & $\times$ \\
离线 RL（Offline RL）& $\times$    & $\checkmark$ & $\checkmark$ & $\times$ \\
\textbf{ETO} & $\checkmark$ & $\checkmark$ & $\checkmark$ & $\checkmark$ \\
\bottomrule
\end{tabular}
\end{table}

只有 ETO 把四样占全了：\textbf{能利用探索失败、由奖励驱动、计算高效、且训练鲁棒}。这也是它的核心卖点——\textbf{既有 RL 的“探索 + 奖励驱动”，又有离线方法的“高效 + 稳定”}。

% =====================================================================
\section{相关工作（简述）}

\textbf{模仿学习}：智能体通过模仿专家示范学策略，BC 是其代表；为缓解 BC 的局限有 DAgger（Ross et al., 2011）等。ETO 与 DAgger\textbf{精神相近}（都从失败案例学），但不同点是：DAgger 会\textbf{在智能体失败处再去收集额外专家轨迹}，而 ETO 是\textbf{从“成败对比对”里学}，不需要为失败处补标专家数据。\textbf{LLM 智能体}：从 AutoGPT、BabyAGI 到 FireAct、AgentTuning、Lumos 等用 GPT-4 造专家数据再 BC 开源模型。\textbf{LLM 策略学习}：从偏好中学（如 RLHF）能提升模型；RLHF 用 PPO 但效率低、不稳，DPO 把目标改写成直接建模偏好；ReST、Math-Shepherd 等也探索了迭代/步级奖励建模。

% =====================================================================
\section{结论、局限与对“AI 应用开发”的启示}

\subsection{核心 takeaways}

\begin{enumerate}[leftmargin=2em,itemsep=4pt]
\item \textbf{“从失败中学”是真有用的}：只学成功（SFT/RFT）会留下性能上限；把\textbf{失败轨迹与成功轨迹配成对}做对比学习，能显著提升智能体表现和泛化能力（OOD 上提升尤其明显）。
\item \textbf{用 DPO 代替在线 RL 是聪明的工程取舍}：DPO 把 RL 的奖励最大化目标改写成“成败偏好对”的对比损失，\textbf{避开了 PPO 的不稳定与低效}，却拿到了“探索 + 奖励驱动”的好处。在 ALFWorld 这种只有二元奖励的环境里，这个优势尤其突出。
\item \textbf{先模仿打底、再试错精炼}：BC 提供可靠起点，ETO 在其上迭代加强。没有 BC 打底时（自我对弈），要先用 RFT 让模型“能跑出成功案例”，再上 ETO。
\item \textbf{迭代有甜区、别贪多}：在固定专家集 + 同一训练集上探索，迭代 2 轮左右最好，超过会因数据多样性不足而\textbf{过拟合}退化。
\item \textbf{粒度选择：整条轨迹对比 > 单步对比}：因为只用最终奖励难以准确判断“单步动作”的好坏（信用分配难题），步级对比反而不稳。
\item \textbf{基础语言能力 $\neq$ 智能体能力}：更强的基座 LLM 不一定造出更强的智能体（Mistral-7B 反不如 Llama-2-7B）。
\end{enumerate}

\subsection{对 AI 应用开发的实际启示}

\begin{itemize}[leftmargin=2em,itemsep=4pt]
\item \textbf{别只喂“正确答案”给你的 Agent}：如果你在用开源模型搭业务智能体（网页操作、工具调用、流程自动化），\textbf{把它线上/测试时跑挂的轨迹收集起来，和成功轨迹配成对做 DPO}，往往比只堆更多“标准示范”更划算、提升更稳。
\item \textbf{优先考虑离线偏好优化（DPO）而非在线 RL}：在线 RL（PPO）调参痛苦、容易崩，尤其当你的环境只能给“成/败”这种粗奖励时。\textbf{DPO + 成败对比对}是更工程友好、更稳的替代方案。
\item \textbf{“失败数据”是被低估的资产}：日志里那些失败 case 不是垃圾——\textbf{它们是免费的、最贴近你真实分布的训练信号}。建立一条“收集失败 → 配对成功 → 周期性 DPO 微调”的流水线，可以让智能体随使用越用越好。
\item \textbf{冷启动策略}：手上没有专家轨迹时，先用拒绝采样（RFT，只保留自己跑成功的轨迹再训）做冷启动，再叠加失败对比学习。
\item \textbf{别盲目迭代}：周期性再训练有“甜区”，盯住验证集，过拟合退化时及时停。
\end{itemize}

\subsection{作者坦承的局限}

\begin{enumerate}[leftmargin=2em,itemsep=3pt]
\item \textbf{信用分配过于简化}：ETO 假设智能体“从一开始就走错了”，于是整条失败轨迹对整条成功轨迹比。但现实中智能体往往是\textbf{中途某一步}才走错；理想做法是定位到那一坏步、只为后续补专家轨迹。可惜\textbf{多数环境不提供这种过程级信息}，难以做步级/过程级奖励建模。潜在方案：用 GPT-4 来识别坏动作、构造细粒度对比数据。
\item \textbf{偏专用、未做通用智能体}：本工作主要针对\textbf{单一具体任务}训练专用智能体，对“强通用智能体”探索不多。未来会研究 ETO 训出的策略的\textbf{可迁移性}、以及多任务训练。
\end{enumerate}

% =====================================================================
\section*{附录要点（精简）}
\addcontentsline{toc}{section}{附录要点}

\textbf{附录 A（数据集细节）}：WebShop 在选“buy”后按属性/价格匹配给奖励；ScienceWorld 含 10 类任务、按子目标给奖励，测试集是“烧开水→熔化铅”这类未见变体；ALFWorld 是 ALFRED 的文字版，分 seen/unseen。CoT 标注：WebShop/ALFWorld 有少量人工轨迹，并用 GPT-4 当教师智能体（取奖励>0.7 的轨迹）；ScienceWorld 用启发式搜索生成黄金轨迹，再用 GPT-4 补 CoT。

\textbf{附录 B（成功率，表 7）}：除平均奖励外另报成功率，ETO 同样全面领先（如 WebShop 37.5、ScienceWorld-Seen 80.3）。

\textbf{附录 C（步级对比细节）}：对前 $t{-}1$ 步用 teacher forcing 喂专家动作，让智能体只预测第 $t$ 步，再按最终奖励把第 $t$ 步的 $(a_t, \hat{a}_t)$ 构造成“好—坏动作对” $a_w\succ a_l\mid u, e^{(t-1)}$，套进 DPO（公式 13）。

\textbf{附录 D（案例研究，图 5/6）}：WebShop 里 SFT 智能体漏选了“3pc”颜色选项导致只拿 0.5 分，ETO 学会正确选择属性拿满分；ScienceWorld 里 SFT 智能体执着于错误动作（盯着“手指画”这种非生物），且动作和自己的 CoT 矛盾，而 ETO 正确找到“青蛙卵→普通蟾蜍”这一动物完成任务。

\textbf{附录 E（评测提示词，图 7/8/9）}：给出 WebShop、ScienceWorld、ALFWorld 三个环境的指令提示词模板（动作格式、可用动作列表、\texttt{Thought/Action} 输出格式等）。

\end{document}
